\section{Introduction}

Visual document retrieval has emerged as a compelling alternative to OCR-based pipelines by directly encoding document page screenshots with vision-language models \citep{faysse2024colpali}. Models such as ColPali and ColQwen represent page images as multi-vector embeddings and perform late-interaction matching \citep{khattab2020colbert,santhanam-etal-2022-colbertv2}, achieving strong results on page-level benchmarks such as ViDoRe.

However, existing visual document retrieval remains largely \textbf{page-level}. While page-level retrieval suits retrieval-augmented generation \citep{lewis2020retrieval}, many practical scenarios, such as enterprise search, academic retrieval, and legal review, require retrieving \textbf{complete documents}. Yet cross-document \textbf{document-level} visual retrieval remains underexplored. ViDoRe \citep{faysse2024colpali,mace2025vidore2} and MMDocIR \citep{dong2025mmdocir} focus on page-level retrieval, while VisDoMBench \citep{suri2025visdombench} targets QA rather than retrieval.

A natural fallback is indexing pages independently and aggregating scores to the document level (e.g., MaxP). This faces a \textbf{dual challenge}: index size scales linearly with total pages, and relevance distributed across multiple pages may not be captured by a single-page score.

We propose \textbf{DocPC}, a document-level visual retrieval framework built upon page-level retrievers via supervised fine-tuning. Its core idea, \textbf{Representative Page Composition (PC)}, selects $K$ representative pages and arranges them into a single grid image, reducing document-side indexing from $\mathcal{O}(N)$ to $\mathcal{O}(1)$ without additional model complexity.

Document-level retrieval also amplifies the \textbf{training challenge} of multi-positive supervision: topically related documents are more common, making standard InfoNCE \citep{oord2018infonce} prone to treating co-relevant documents as negatives. We combine \textbf{Multi-Positive Contrastive Loss} \citep{khosla2020supervised,rivas2024conked} with a sparsely scheduled ApproxNDCG-based listwise loss \citep{qin2010general} to address this.

We further construct \textbf{DocViRe}, a benchmark for document-level visual retrieval with multi-positive relevance annotations across seven domains. Our contributions:
\begin{itemize}
    \item We formalize document-level visual retrieval and propose Representative Page Composition for $\mathcal{O}(1)$ document indexing.
    \item We construct DocViRe, a dedicated document-level visual retrieval  benchmark with multi-positive annotations.
    \item We empirically validate DocPC on DocViRe, demonstrating consistent retrieval gains from page selection and multi-positive training while significantly reducing index size.
\end{itemize}